% spellchecker:disable
%pagenumbering{null}

\

\cleardoublepage

\

\pagestyle{empty}

\IfLanguageName{ngerman}{
\textbf{Eidesstattliche Versicherung}\\

Hiermit versichere ich an Eides Statt,
dass ich diese Arbeit bzw. im Fall einer Gruppenarbeit
den von mir entsprechend gekennzeichneten Anteil an der Arbeit
selbständig verfasst habe.
Ich habe keine unzulässige Hilfe Dritter in Anspruch genommen.
Zudem habe ich keine anderen
als die angegebenen Quellen und Hilfsmittel
benutzt und alle Ausführungen (insbesondere Zitate),
die anderen Quellen wörtlich oder sinngemäß entnommen wurden,
kenntlich gemacht.
Die Liste zugelassener Hilfmittel vom \DTMDate{\dateToolsAndAI} ist mir ausdrücklich bekannt.

Ich versichere,
dass die von mir in elektronischer Form eingereichte Version dieser Arbeit
mit den eingereichten gedruckten Exemplaren übereinstimmt.

Mir ist bekannt, dass im Falle eines Täuschungsversuches
die betreffende Leistung als mit “nicht ausreichend” (5,0) bewertet gilt.
Zudem kann ein Täuschungsversuch als Ordnungswidrigkeit
mit einer Geldbuße von bis zu 50.000 Euro geahndet werden.
Im Falle eines mehrfachen oder sonstigen schwerwiegenden Täuschungsversuchs
kann ich zudem exmatrikuliert werden.

Mir ist bekannt,
dass sich die Prüferin oder der Prüfer bzw. der Prüfungsausschuss
zur Feststellung der Täuschung
des Einsatzes einer entsprechenden Software
oder sonstiger elektronischer Hilfsmittel bedienen kann.

Ich versichere an Eides Statt,
dass ich die vorbenannten Angaben nach bestem Wissen und Gewissen gemacht habe
und dass die Angaben der Wahrheit entsprechen
und ich nichts verschwiegen habe.

Die Strafbarkeit einer falschen eidesstattlichen Versicherung ist mir bekannt,
insbesondere die Strafandrohung gemäß §§ 156, 161 StGB,
auf welche ich konkret hingewiesen wurde.

\begin{tcolorbox}[colframe=black!100!white,colback=black!0!white,sharp corners]
\footnotesize
\textbf{§ 156 Falsche Versicherung an Eides Statt}\\
Wer vor einer zur Abnahme einer Versicherung an Eides Statt zuständigen Behörde
eine solche Versicherung falsch abgibt
oder unter Berufung auf eine solche Versicherung falsch aussagt,
wird mit Freiheitsstrafe bis zu drei Jahren oder mit Geldstrafe bestraft.

\tcblower

\footnotesize
\textbf{§ 161 Fahrlässiger Falscheid; fahrlässige falsche Versicherung an Eides Statt}\\
(1) Wenn eine der in den §§ 154 bis 156 bezeichneten Handlungen aus Fahrlässigkeit begangen worden ist,
so tritt Freiheitsstrafe bis zu einem Jahr oder Geldstrafe ein.
(2) Straflosigkeit tritt ein,
wenn der Täter die falsche Angabe rechtzeitig berichtigt.
Die Vorschriften des § 158 Abs. 2 und 3 gelten entsprechend.
\end{tcolorbox}

\vfill
Duisburg, \today\\
$\overline{\parbox{4.8cm}{(Ort, Datum)}} ~~~~~~~~~~~~~~~~~~~~~~~~~~~ \overline{\parbox{7cm}{(Unterschrift)}}$
}
{
\textbf{Declaration in lieu of an oath}\\

I hereby solemnly declare that I have independently completed this work or,
in the case of group work,
the part of the work that I have marked accordingly.
I have not made use of the unauthorised assistance of third parties.
Furthermore, I have used only the stated sources or aids
and I have referenced all statements (particularly quotations)
that I have adopted from the sources I have used verbatim or in essence.
The list of allowed tools from \DTMDate{\dateToolsAndAI} is explicitly known to me.

I declare that the version of the work I have submitted in digital form
is identical to the printed copies submitted.

I am aware that, in the case of an examination offence,
the relevant assessment will be marked as ‘insufficient’ (5.0).
In addition, an examination offence
may be punishable as an administrative offence (Ordnungswidrigkeit)
with a fine of up to €50,000.
In cases of multiple or otherwise serious examination offences,
I may also be removed from the register of students.

I am aware that the examiner and/or the Examination Board
may use relevant software or other electronic aids
in order to establish an examination offence has occurred.

I solemnly declare that I have made the previous statements to the best of my knowledge and
belief and that these statements are true and I have not concealed anything.

I am aware of the potential punishments for a false declaration in lieu of oath
and in particular of the penalties
set out in Sections 156 and 161 of the German Criminal Code (Strafgesetzbuch; StGB),
which I have been specifically referred to.

\begin{tcolorbox}[colframe=black!100!white,colback=black!0!white,sharp corners]
\footnotesize
\textbf{Section 156 False declaration in lieu of an oath}

Whoever falsely makes a declaration in lieu of an oath
before an authority which is competent to administer such declarations
or falsely testifies whilst referring to such a declaration
incurs a penalty of imprisonment for a term not exceeding three years or a fine.

\tcblower

\footnotesize
\textbf{Section 161 Negligent false oath; negligent false declaration in lieu of oath}

(1) Whoever commits one of the offences referred to in Sections 154 to 156 by negligence
incurs a penalty of imprisonment for a term not exceeding one year or a fine.
(2) No penalty is incurred if the offender corrects the false statement in time.
The provisions of Section 158 (2) and (3) apply accordingly.
\end{tcolorbox}

\vfill
Duisburg, \today\\
$\overline{\parbox{4.8cm}{(Place, Date)}} ~~~~~~~~~~~~~~~~~~~~~~~~~~~ \overline{\parbox{7cm}{(Signature)}}$
}
